\documentclass{resume} % Use the custom resume.cls style
\usepackage[left=0.4in,top=0.4in,right=0.4in,bottom=0.4in]{geometry} % Document margins
\usepackage[usenames,dvipsnames]{xcolor}
\hypersetup{colorlinks=true,linkcolor=MidnightBlue,urlcolor=MidnightBlue} % Configures hyperref since it is already loaded by resume.cls
\newcommand{\tab}[1]{\hspace{.2667\textwidth}\rlap{#1}}
\newcommand{\itab}[1]{\hspace{0em}\rlap{#1}}
\name{HEITOR GOMES DIAS JASKO} % Your name
\address{Palmares Paulista, SP | +55 (17) 99670-2090}
\address{\href{https://www.linkedin.com/in/heitor-gomes-dias-jasko-9694b8273}{linkedin} \href{mailto:heitor@jasko.dev}{heitor@jasko.dev} \ \href{https://www.jasko.dev}{portfolio}} %
\begin{document}
%----------------------------------------------------------------------------------------
%	OBJECTIVE
%----------------------------------------------------------------------------------------
\begin{rSection}{OBJETIVO}
Estágio em Desenvolvimento de Software. Backend em \textbf{Java/Spring Boot} e frontend em \textbf{React/Next.js}, com vivência em infraestrutura própria (\textbf{Docker, CI/CD, cloud}) e integração de \textbf{IA (RAG)} a produtos.
\end{rSection}
%----------------------------------------------------------------------------------------
%	EDUCATION SECTION
%----------------------------------------------------------------------------------------
\begin{rSection}{FORMAÇÃO}
\textbf{Análise e Desenvolvimento de Sistemas}, Instituto Federal de São Paulo \hfill {Conclusão: 2028}\\
\textbf{Técnico em Desenvolvimento de Sistemas}, Etec Elias Nechar \hfill {Conclusão: 2024}
\end{rSection}
%----------------------------------------------------------------------------------------
% TECHINICAL STRENGTHS 
%----------------------------------------------------------------------------------------
\begin{rSection}{HABILIDADES}
\begin{tabular}{ @{} >{\bfseries}l @{\hspace{6ex}} p{5.8in} }
Desenvolvimento & Java (Spring Boot, JPA, Security), API REST, React, Next.js, TypeScript \\
Banco de Dados & PostgreSQL, SQL, Modelagem de Dados \\
IA Aplicada & Integração de LLMs, RAG (Retrieval-Augmented Generation), Prompt Engineering \\
DevOps \& CI/CD & Docker, Nginx, GitHub Actions, Pipelines de Deploy Automatizado, Git/GitHub \\
Infraestrutura & Linux, Oracle Cloud (OCI), Cloudflare, Portainer, Netdata \\
Ferramentas & VS Code, Postman, Google Workspace, Office \\
Idiomas & Inglês (Intermediário) \\
\end{tabular}
\end{rSection}
%----------------------------------------------------------------------------------------
%	WORK EXPERIENCE SECTION
%----------------------------------------------------------------------------------------
\begin{rSection}{EXPERIÊNCIA}
\textbf{Suporte técnico em software de gerenciamento empresarial} \hfill Até março, 2026\\
Solução Sistemas \hfill \textit{Catanduva, SP}
 \begin{itemize}
    \itemsep -5pt {} 
    \item Diagnostiquei erros e \textbf{corrigi} inconsistências em bancos de dados por meio de consultas \textbf{SQL}, além de extrair relatórios para clientes.
    \item Gerenciei atualizações no ambiente do cliente, garantindo estabilidade e continuidade das novas versões.
    \item Solucionei problemas (troubleshooting) em processos fiscais e regras de negócio.
    \item Validei processos de compilação e build, assegurando a integridade dos executáveis antes da distribuição.
    \item \textbf{Identifiquei} documentação técnica dispersa em arquivos soltos em servidor SMB, \textbf{otimizei e migrei} o conteúdo markdown para uma documentação em Vite (React); \textbf{implementei} um assistente de \textbf{IA com RAG} para consulta direta pela equipe.
\end{itemize}
\textbf{Suporte técnico em software de gerenciamento empresarial} \hfill Até setembro, 2025\\
Visualsoft Brasil \hfill \textit{Catanduva, SP}
 \begin{itemize}
    \itemsep -5pt {} 
    \item Diagnostiquei e resolvi incidentes técnicos via helpdesk, cumprindo prazos de atendimento (\textbf{SLA}).
    \item Analisei erros na emissão de \textbf{NF-e/NFC-e} (XML) e dei suporte a rotinas de banco de dados e controle de estoque.
    \item \textbf{Desenvolvi} um script em \textbf{Python} que automatizou a instalação e ativação do sistema legado, processo antes manual e repetitivo, reduzindo-o a uma única execução.
 \end{itemize}
\end{rSection}
%----------------------------------------------------------------------------------------
%	PROJECTS SECTION
%----------------------------------------------------------------------------------------
\begin{rSection}{PROJETOS}
\begin{itemize}
    \itemsep -5pt {}
    \item \textbf{Benius} (\href{https://benius.jasko.dev}{benius.jasko.dev}) — Rede social anônima estilo Twitter; \textbf{desenvolvi} backend \textbf{Java}/\textbf{PostgreSQL} e frontend em \textbf{Next.js/React}.
    \item \textbf{Tasty} (\href{https://tasty.jasko.dev}{tasty.jasko.dev}) — Simulador de terminal web em Java, com autenticação segura e chat assíncrono integrado a IA.
    \item \textbf{Infraestrutura própria} — \textbf{Projetei} e mantenho servidor Oracle Cloud (OCI) com Docker, Nginx e Cloudflare hospedando as aplicações acima; \textbf{implementei CI/CD} via GitHub Actions (deploy automático a cada push na main, estilo Vercel).
\end{itemize}
\end{rSection}
\end{document}